\documentclass{article}
\usepackage{hyperref}
\usepackage[margin=0.5in]{geometry}
\usepackage{amsmath}
\usepackage{graphicx}
\graphicspath{ {./assets/} }

\title{PHYSIOL 578: HW 1.3\footnotetext{U-M GPT was used to refine language of answers derived independently.}}
\author{Jack Yang \\ \href{mailto:yhyunjoo@umich.edu}{yhyunjoo@umich.edu}}
\date{September 27, 2026}

\begin{document}
\maketitle
\section{
    Identify a ligand-gated ion channel that is not in the course materials.
    Provide a full description of the receptor, its ligands, and any clinical targets.
    Relate the receptor function to human physiology. (1 page max)
}
The glycine receptor (GlyR) is a member of the pentameric Cys-loop family of ligand-gated ion channels.
It is assembled from five pore-forming subunits, typically $\alpha$ and $\beta$ subunits.
Adult receptors are commonly $\alpha$1$\beta$ heteromers, while the $\beta$ subunit helps anchor the receptor to the postsynaptic membrane through the scaffolding protein gephyrin.

Glycine is the primary endogenous agonist.
Other endogenous amino acids, including $\beta$-alanine and taurine, can act as partial agonists.
Strychnine is a potent competitive antagonist, whereas compounds such as picrotoxin can inhibit GlyRs through noncompetitive channel block.
Ethanol, cannabinoids, and ivermectin can positively modulate certain GlyR subtypes, although their effects depend on receptor composition and concentration.

When activated, GlyRs conduct primarily $Cl^-$.
Opening the channel drives the membrane potential toward the chloride equilibrium potential and increases membrane conductance.
In mature neurons, this usually produces hyperpolarizing and/or shunting inhibition, reducing the likelihood of action-potential firing.
GlyRs are especially abundant in the spinal cord, brainstem, and retina, where they regulate motor activity, sensory processing, reflexes, breathing, and pain signaling.

Mutations affecting GlyR subunits—particularly \textit{GLRA1} and \textit{GLRB}—or the glycine transporter GlyT2 can cause hyperekplexia.
This disorder is characterized by an exaggerated startle response followed by transient generalized muscle stiffness, which may cause unprotected falls or breathing difficulties in infants.
Clonazepam is commonly used to treat hyperekplexia by strengthening GABAergic inhibition, although it does not directly activate GlyRs.
GlyR-positive allosteric modulators, particularly those targeting the $\alpha$3 subtype involved in spinal pain pathways, are also being investigated as potential analgesic therapies.

\section{
    Provide a definition of excitability (your choice).
    Use the \href{gaba.tonydefazio.com}{gaba.tonydefazio.com} website to measure excitability according to your definition using the cortical neuron model and the GnRH neuron model.
    Predict the change in excitability in the two models if you were to change one of the parameters in the simulation (GABA PSC frequency, Rin, etc).
    Quantify your observations. (1 page max)
}
We operationally define neuronal excitability as the reciprocal of the \textbf{Rheobase} ($I_{\text{rheo}}$)—the minimum depolarizing current step required to elicit at least one action potential from rest:
$$\text{Excitability} = \frac{1}{\text{Rheobase}} \quad (\text{pA}^{-1})$$
Under this definition, a neuron requiring a smaller current step to reach spike threshold exhibits \textbf{higher excitability} (larger $\text{pA}^{-1}$ value), whereas a neuron requiring a larger current step exhibits \textbf{lower excitability} (smaller $\text{pA}^{-1}$ value).

Using the default current-clamp simulation protocols on \href{https://gaba.tonydefazio.com}{gaba.tonydefazio.com}, the baseline \textbf{cortical pyramidal neuron model} exhibits a rheobase of $I_{\text{rheo}} = 600\text{ pA}$, yielding an excitability of:
$$\text{Excitability}_{\text{cortical}} = \frac{1}{600\text{ pA}} = 1.67 \times 10^{-3}\text{ pA}^{-1} \quad (0.00167\text{ pA}^{-1})$$
In contrast, the baseline \textbf{mouse GnRH neuron model} exhibits a rheobase of $I_{\text{rheo}} = 10\text{ pA}$, yielding an excitability of:
$$\text{Excitability}_{\text{GnRH}} = \frac{1}{10\text{ pA}} = 1.00 \times 10^{-1}\text{ pA}^{-1} \quad (0.100\text{ pA}^{-1})$$
At baseline, the GnRH neuron is \textbf{60.0 times more excitable} than the cortical pyramidal cell ($\frac{0.100}{0.00167} = 60$). This substantial difference arises because GnRH neurons possess a higher intrinsic input resistance ($R_{in}$) and smaller membrane capacitance ($C_m$), generating significantly larger voltage deflections ($\Delta V = I \cdot R_{in}$) for a given current injection.

\textbf{Parameter Shift Prediction \& Observation}: We predict that shifting $E_{\text{GABA}}$ to a more depolarized potential ($0\text{ mV}$) will increase inward synaptic driving force during active GABA PSCs, thereby decreasing the Rheobase and enhancing excitability in both models.
\begin{itemize}
    \item \textbf{Cortical Pyramidal Neuron}: Shifting $E_{\text{GABA}}$ from $-80\text{ mV}$ to $0\text{ mV}$ reduces $I_{\text{rheo}}$ from $600\text{ pA}$ to $400\text{ pA}$, increasing excitability to $2.50 \times 10^{-3}\text{ pA}^{-1}$ (\textbf{+50.0\% increase}).
    \item \textbf{GnRH Neuron}: Shifting $E_{\text{GABA}}$ from $-36.5\text{ mV}$ to $0\text{ mV}$ reduces $I_{\text{rheo}}$ from $10\text{ pA}$ to $2\text{ pA}$, increasing excitability to $5.00 \times 10^{-1}\text{ pA}^{-1}$ (\textbf{+400.0\% increase}).
\end{itemize}

\begin{table}[h!]
\centering
\begin{tabular}{|l|c|c|c|c|}
\hline
\textbf{Model \& Condition} & $E_{\text{GABA}}$ \textbf{(mV)} & \textbf{Rheobase (pA)} & \textbf{Excitability ($\text{pA}^{-1}$)} & \textbf{Relative Shift} \\ \hline
Cortical (Baseline) & $-80.0\text{ mV}$ & $600\text{ pA}$ & $1.67 \times 10^{-3}\text{ pA}^{-1}$ & Baseline \\ \hline
Cortical ($E_{\text{GABA}} = 0\text{ mV}$) & $0.0\text{ mV}$ & $400\text{ pA}$ & $2.50 \times 10^{-3}\text{ pA}^{-1}$ & \textbf{+50.0\% Increase} \\ \hline
GnRH (Baseline) & $-36.5\text{ mV}$ & $10\text{ pA}$ & $1.00 \times 10^{-1}\text{ pA}^{-1}$ & Baseline \\ \hline
GnRH ($E_{\text{GABA}} = 0\text{ mV}$) & $0.0\text{ mV}$ & $2\text{ pA}$ & $5.00 \times 10^{-1}\text{ pA}^{-1}$ & \textbf{+400.0\% Increase} \\ \hline
\end{tabular}
\end{table}
\end{document}